% UCL Impact Statement — required for all UCL PhD theses.
% Limit: 500 words (UCL regulations).
% DRAFT 13 Sep 2026 for Damian's review; dissemination paragraph needs
% his details.

This thesis measures two features of English secondary schools that matter
for pupils but have been hard to observe: the culture a school actually
practises day to day, and whether it offers a subject that leads to
well-paid careers. Both sets of findings
bear directly on live policy questions.

The findings on school culture speak first to how schools are monitored.
England is reforming school inspection, and one question in that reform is
what an inspection report can and cannot tell parents. This thesis provides
an answer grounded in measurement. The text of an inspection report gives a
modest indication of how warm a school is, and a modest indication of how
strict it is. The two indications come from different parts of the report.
Schools' own documents reflect what headteachers believe rather than what
visitors observe. The methods developed here use language models to read
inspection reports, behaviour policies and websites. Their scores were
checked against ratings made by observers on school visits. They could be
reused by researchers or by government to track features of schooling that
current data do not capture, at national scale and at low cost.

The findings speak second to the public debate about school behaviour
policy. That debate is often conducted as a contest between strict
discipline and pupil wellbeing. The evidence here is that warmth and
strictness are separate features of a school. A school can be high on one
and low on the other. Pupils make the most academic progress at schools
that have both. Warmth predicts pupil progress as strongly as strictness
does. This holds only when warmth is measured by visiting observers rather
than by the headteacher's own description of the school. For school
leaders, one finding is a practical caution. Headteachers describe their
schools as warmer than observers find them. A headteacher's own sense of
how warm the school is may therefore be the least reliable part of their
picture of it.

The third paper's impact is the most direct. During the period the paper
studies, fewer than half of English state schools offered A-level
economics, the principal route into one of the
highest-earning degrees. Whether a school offers the subject usually comes
down to one thing. It needs at least one teacher qualified to teach
economics. The paper shows that hiring an economics teacher more than pays
for itself. Pupils who take the subject go on to earn more and pay more
tax. On the paper's main estimate, that extra tax exceeds the full cost of
recruiting, training and employing the teacher several times over. Economics currently receives no teacher-training bursary, unlike
other shortage subjects; the paper's calculations provide a costed case
for one. Because the schools without provision are concentrated outside London and
the South East, expanding provision would also broaden regional access to
the subject and to the careers it leads to. These results are of direct
relevance to the Department for Education, to teacher-supply policy, and
to the Bank of England and Royal Economic Society campaigns to widen
participation in economics.

[Dissemination paragraph to complete: seminar and conference
presentations, engagement with schools and participating headteachers,
co-authors' policy networks (IFS), and any working-paper or publication
plans.]
